\subsection{Example: de Rham complex}

In this section, we assume that A is a commutative \(k\) -algebra over a commutative ring \(k\) . One denotes by \(\Omega_{\mathbf{A}/k}^{1}\) the A-module of Kähler differentials of A (III, p.~134), \(d^{0}: \mathbf{A} \to \Omega_{\mathbf{A}/k}^{1}\) the \(k\) -derivation \(d_{\mathbf{A}/k}\) , and \(\Omega_{\mathbf{A}/k}\) the graded \(k\) -algebra \(\boldsymbol{\Lambda}_{\mathbf{A}}(\Omega_{\mathbf{A}/k}^{1})\) .

PROPOSITION 13. — There exists a unique k-antiderivation \(d: \Omega_{\mathbf{A}/k} \to \Omega_{\mathbf{A}/k}\) of degree 1, with square zero, extending the derivation \(d^{0}: \mathbf{A} \to \Omega_{\mathbf{A}/k}^{1}\) .

Let us prove the uniqueness of the antiderivation \(d\) . Since \(d \circ d = 0\) we have, for \(y, x_1, \ldots, x_p \in \mathbf{A}\) :

\[
d \left(y d x _ {1} \wedge \dots \wedge d x _ {p}\right) = d y \wedge d x _ {1} \wedge \dots \wedge d x _ {p}.
\]

The A-module \(\Omega_{\mathbf{A}/k}^{p}\) being generated by the elements \(dx_{1} \wedge \ldots \wedge dx_{p}\) , this proves the uniqueness of \(d\) .

To prove existence, it suffices to construct a \(k\) -homomorphism \(d^{1}:\Omega_{\mathbb{A}/k}^{1}\to\Omega_{\mathbb{A}/k}^{2}\) such that \(d^{1}\circ d^{0}=0\) and

\[
d ^ {1} (a \omega) = d ^ {0} (a) \wedge \omega + a d ^ {1} (\omega) \quad \text {for} \quad a \in A, \quad \omega \in \Omega_ {A k} ^ {1}.\tag{11}
\]

Indeed, it then follows from III, p.~128, prop. 14 (taking into account III, p.~118, remark 2) that there exists an antiderivation \(d: \Omega_{\mathbf{A}/k} \to \Omega_{\mathbf{A}/k}\) that coincides with \(d^{0}\) in degree 0 and with \(d^{1}\) in degree 1. Since \(d^{0}\) vanishes on \(\mathbf{A}\) , the antiderivation \(d\) is \(k\) -linear; since \(d^{1} \circ d^{0} = 0\) , we have \(d \circ d = 0\) since \(\Omega_{\mathbf{A}/k}\) is generated as \(\mathbf{A}\) -algebra by the elements \(d^{0} a\) for \(a \in \mathbf{A}\) .

To define \(d^{1}\) , recall (III, p. 133) that \(\Omega_{\mathbf{A}/k}^{1}\) is equal to the A-module \(\mathfrak{I}/\mathfrak{I}^{2}\) , where \(\mathfrak{I}\) is the kernel of the multiplication \(m: \mathbf{A} \otimes_{k} \mathbf{A} \to \mathbf{A}\) . Consider the \(k\) -linear map \(u: \mathbf{A} \otimes_{k} \mathbf{A} \to \Omega_{\mathbf{A}/k}^{2}\) defined by \(u(x \otimes y) = d^{0}(y) \wedge d^{0}(x)\) . We have

\[
u (a x \otimes y - x \otimes a y) = d ^ {0} (y) \wedge d ^ {0} (a x) - d ^ {0} (a y) \wedge d ^ {0} (x) = d ^ {0} (x y) \wedge d ^ {0} (a)
\]

for x, y, and a in A, whence

\[
u ((a \otimes 1 - 1 \otimes a) \xi) = d ^ {0} (m (\xi)) \wedge d ^ {0} (a), \quad \xi \in \mathrm{A} \otimes_ {k} \mathrm{A}, \quad a \in \mathrm{A}. \tag {12}
\]

Since \(\mathfrak{I}\) is generated as a left A-module by the elements ( \(a \otimes 1 - 1 \otimes a\) ) for \(a \in \mathbf{A}\) , we deduce that \(u(\mathfrak{I}^2) = 0\) ; therefore \(u\) defines by restriction to \(\mathfrak{I}\) and passage to the quotient a \(k\) -linear map \(d^1: \mathfrak{I}/\mathfrak{I}^2 \to \Omega_{\mathrm{A}/k}^2\) .

By making \(\xi = b \otimes 1\) in (12), with \(b \in \mathbf{A}\) , we obtain \(d^{1}(bd^{0}(a)) = d^{0}(b) \wedge d^{0}(a)\) ; it follows that \(d^{1} \circ d^{0} = 0\) and that \(d^{1}(c\omega) = d^{0}(c) \wedge \omega + cd^{1}(\omega)\) for \(c \in \mathbf{A}\) and \(\omega = bd^{0}(a)\) . Since \(\Omega_{\mathbf{A}/k}^{1}\) is generated as \(k\) -module by the elements \(bd^{0}(a)\) , for \(a\) and \(b\) in \(\mathbf{A}\) , formula (11) is satisfied for all \(\omega \in \Omega_{\mathbf{A}/k}\) , which completes the proof of the proposition.

Sometimes it is said that the elements \(\omega\in\Omega_{A/k}^{p}\) are the exterior differential forms of degree \(p\) of A over \(k\) and that the antiderivation \(d\) is the exterior differential of \(\Omega_{A/k}\) ; the complex ( \(\Omega_{A/k}\) , \(d\) ) is called the de Rham complex of A over \(k\) , and its homology is the de Rham cohomology of A over \(k\) .

Example 1. --- Take for A the ring \(k[\mathbf{X}_{1},\ldots,\mathbf{X}_{n}]\) . Then \(\Omega_{\mathbf{A}/k}^{1}\) is a free A-module with basis \(d\mathbf{X}_{1},\ldots,d\mathbf{X}_{n}\) (III, p.~134, example). Consequently, if for every subset \(\mathrm{I}=\{i_{1},\ldots,i_{p}\}\) of \([1,n]\) we set \(d\mathbf{X}_{\mathrm{I}}=d\mathbf{X}_{i_{1}}\wedge\ldots\wedge d\mathbf{X}_{i_{p}}\) (with \(i_{1}<\ldots<i_{p}\) ), the A-module \(\Omega_{\mathbf{A}/k}^{p}\) admits as a basis the elements \(d\mathbf{X}_{\mathrm{I}}\) , where I ranges over the set of subsets of \([1,n]\) of cardinality \(p\) . We have

\[
d (\mathrm{P} d \mathrm{X} _ {\mathrm{I}}) = d \mathrm{P} \wedge d \mathrm{X} _ {\mathrm{I}} = \sum_ {i \notin \mathrm{I}} (- 1) ^ {n (\mathrm{I}, i)} \frac {\partial \mathrm{P}}{\partial \mathrm{X} _ {i}} d \mathrm{X} _ {\mathrm{I} \cup \{i \}},
\]

where \(n(\mathbf{I}, i)\) denotes the number of elements of I strictly less than \(i\) .

The cycles of \(Z^{p}(\Omega_{\mathbf{A}/k})\) are therefore the elements \(\omega = \sum_{\text{Card (I)} = p} \mathrm{P}_{\mathrm{I}} d\mathrm{X}_{\mathrm{I}}\) such that, for every subset \(J\) with \((p + 1)\) elements of \([1, n]\) , one has:

\[
\sum_ {i \in J} (- 1) ^ {n (J, i)} \frac {\partial P _ {J - \{i \}}}{\partial X _ {i}} = 0.
\]

The element ω is a boundary if one can choose, for every subset \(J \subset [1, n]\) with \((p - 1)\) elements, a polynomial \(Q_{J} \in A\) such that:

\[
\mathrm{P} _ {\mathrm{I}} = \sum_ {j \in \mathrm{I}} (- 1) ^ {n (\mathrm{I}, j)} \frac {\partial \mathrm{Q} _ {\mathrm{I} - \{j \}}}{\partial \mathrm{X} _ {j}}.
\]

We will see in § 9 that the de Rham complex of A over k is acyclic in degrees \textgreater{} 0 if k is a Q-algebra (X, p.~159, remark 4).

* Example 2. --- Suppose that \(k = \mathbf{C}\) and \(\mathbf{A} = \mathbf{C}\) \([\mathrm{X}_1, \ldots, \mathrm{X}_n] / (\mathrm{P}_1, \ldots, \mathrm{P}_r)\) , where the \(\mathrm{P}_i\) are polynomials in \(\mathrm{X}_1, \ldots, \mathrm{X}_n\) , such that the set of points of \(\mathbf{C}^n\) where all the \(\mathrm{P}_i\) vanish is an analytic submanifold V of \(\mathbf{C}^n\) . One can show that the de Rham cohomology of A over C is isomorphic to the singular cohomology H(V, C). *

Let M be an A-module and \(\nabla^{0}\) a \(k\) -linear map from M into \(\mathbf{M} \otimes_{\mathbf{A}} \Omega_{\mathbf{A}/k}^{1}\) such that

\[
\nabla^ {0} (a m) = a \nabla^ {0} (m) + m \otimes d a \quad \text { for } \quad a \in \mathrm{A}, \quad m \in \mathrm{M}\tag{13}
\]

(one sometimes says that \(\nabla^{0}\) is a connection on the A-module M).

Proposition 14. --- (i) There exists a unique k-linear map \(\nabla\) of the \(\Omega_{\mathbf{A}/k}\) right module \(\mathbf{M} \otimes_{\mathbf{A}} \Omega_{\mathbf{A}/k}\) into itself, graded of degree 1, which extends \(\nabla^{0}\) in degree 0 and satisfies the identity:

\[
\nabla (x \omega) = (\nabla x) \omega + (- 1) ^ {p} x (d \omega) \quad \mathrm{for} \quad x \in \mathrm{M} \otimes_ {\mathrm{A}} \Omega_ {\mathrm{A} / k} ^ {p}, \quad \omega \in \Omega_ {\mathrm{A} / k}. \tag {14}
\]

\begin{enumerate}
\def\labelenumi{(\roman{enumi})}
\setcounter{enumi}{1}
\tightlist
\item
  The composite mapping \(\nabla \circ \nabla\) is \(\Omega_{\mathbf{A} / k}\) -linear; in particular the map \(\mathbf{R} = \nabla^{1} \circ \nabla^{0}\) of \(\mathbf{M}\) to \(\mathbf{M} \otimes_{\mathbf{A}} \Omega_{\mathbf{A} / k}^{2}\) is \(\mathbf{A}\) -linear, and we have
\end{enumerate}

\[
\nabla \circ \nabla (m \otimes \omega) = R (m). \omega \quad \mathrm{for} \quad m \in M, \quad \omega \in \Omega_ {A / k}.
\]

The homomorphism R is sometimes called the curvature homomorphism of the connection \(\nabla^{0}\) ; if it is zero, the pair \((\mathbf{M} \otimes_{\mathbf{A}} \Omega_{\mathbf{A}/k}, \nabla)\) is a complex, also called the de Rham complex of \((\mathbf{M}, \nabla^{0})\) over \(k\) .

Let us prove (i). The uniqueness of \(\nabla\) is obvious. Let us define a \(k\) -homomorphism \(\overline{\nabla}\) of \(\mathbf{M} \otimes_k \Omega_{\mathbf{A}/k}\) to \(\mathbf{M} \otimes_{\mathbf{A}} \Omega_{\mathbf{A}/k}\) by

\[
\overline {{{\nabla}}} (m \otimes_ {k} \omega) = (\nabla^ {0} m) \omega + m \otimes d \omega \quad \text { for } \quad m \in \mathbf {M}, \quad \omega \in \Omega_ {\mathrm{A} / k}.
\]

It follows from (13) that \(\overline{\nabla}(am \otimes \omega) = \overline{\nabla}(m \otimes a\omega)\) , so that by passing to the quotient we obtain a \(k\) -homomorphism \(\nabla\) of \(\mathbf{M} \otimes_{\mathbf{A}} \Omega_{\mathbf{A}/k}\) into itself, graded of degree 1, extending \(\nabla^{0}\) in degree 0. Let us verify (14): we have, for \(m \in \mathbf{M}\) , \(\alpha \in \Omega_{\mathbf{A}/k}^{p}\) , \(\omega \in \Omega_{\mathbf{A}/k}\) :

\[
\begin{array}{r l} \nabla ((m \otimes \alpha). \omega) & = \nabla (m \otimes (\alpha \wedge \omega)) = \nabla^ {0} (m). (\alpha \wedge \omega) + m \otimes d (\alpha \wedge \omega) \\ & = \nabla^ {0} (m) \alpha . \omega + (m \otimes d \alpha) \omega + (- 1) ^ {p} (m \otimes \alpha) d \omega \\ & = (\nabla (m \otimes \alpha)) \omega + (- 1) ^ {p} (m \otimes \alpha) d \omega \end{array}
\]

which proves (14) for \(x = m \otimes \alpha\) ; the general case follows from this by linearity.

Let us prove (ii). Let \(x \in \mathbf{M} \otimes_{\mathbf{A}} \Omega_{\mathbf{A}/k}^{p}\) , \(\omega \in \Omega_{\mathbf{A}/k}\) ; by repeated application of (14), we obtain:

\[
\begin{array}{r l} \nabla \circ \nabla (x \omega) = & \nabla (\nabla (x) \omega) + (- 1) ^ {p} \nabla (x d \omega) \\ & = (\nabla \circ \nabla (x)) \omega + (- 1) ^ {p + 1} \nabla (x) (d \omega) + (- 1) ^ {p} \nabla (x) (d \omega) \\ & = (\nabla \circ \nabla (x)) \omega , \end{array}
\]

which proves the first assertion of (ii); the others follow immediately.

